This conclusion extends directly to systems containing any number of identical particles.
The number of particles may be arbitrary.
To see this, choose any pair of them.
Suppose, for example, that this pair has the property of being described by symmetric wave functions.
The identity of the particles then requires every other pair of the same kind to have this property as well.
Thus, in one possibility, the wave function remains wholly unchanged when any pair of particles is exchanged.
It consequently remains unchanged under any permutation of the particles whatsoever.
In the other possibility, it changes sign whenever a pair is exchanged.
In the first case the wave function is called \emph{symmetric}.
In the second it is called \emph{antisymmetric}.
Which of these two properties applies depends on the kind of particle.
Particles described by antisymmetric functions are said to obey \emph{Fermi--Dirac statistics}.
They are also called \emph{fermions}.
Particles described by symmetric functions are said to obey \emph{Bose--Einstein statistics}.
They are also called \emph{bosons}\footnote{This terminology comes from the names of the statistics used to describe an ideal gas. In one case the gas consists of particles whose wave functions are antisymmetric. In the other it consists of particles whose wave functions are symmetric. In fact, the distinction here is not merely between two statistics. The mechanics themselves are also essentially different. Fermi statistics was proposed by Fermi (E.~Fermi) for electrons in 1926. Its connection with quantum mechanics was established by Dirac in 1926. Bose statistics was proposed by Bose (S.~Bose) for light quanta. Einstein generalized it in 1924.}.
